\documentclass[11pt]{article}
\usepackage{preamble}
\setlength{\parskip}{4pt}

\begin{document}

\daybanner{AP Statistics \textperiodcentered\ Unit 4 \textperiodcentered\ Topic 4.1 \textperiodcentered\ TEACHER KEY}{Two Production Lines, Two Sampling Plans}

\sectionbanner{CED TETHER}

{\small
\begin{itemize}[leftmargin=1.5em,itemsep=2pt,topsep=2pt]
  \item \textbf{Enduring Understanding: UNC-3} --- Probabilistic reasoning allows us to anticipate patterns in data.
  \item Skills:
  \item \textbf{Skill 3.B} --- Determine parameters for probability distributions.
  \item \textbf{Skill 3.C} --- Describe probability distributions.
  \item \textbf{Skill 4.B} --- Interpret statistical calculations and findings to assign meaning or assess a claim.
  \item Learning Objectives and Essential Knowledge:
  \item UNC-3.Q Determine parameters for a sampling distribution for sample means. [Skill 3.B]
  \item UNC-3.Q.1 For a numerical variable, when random sampling from a population with mean $\mu$ and standard deviation $\sigma$, the sampling distribution of the sample mean has mean $\mu$\_$\bar{x}$ = $\mu$ and standard deviation $\sigma$\_$\bar{x}$ = $\sigma$/$\surd$n. UNC-3.Q.2 If sampling without replacement and sample size is less than 10\% of the population size, the difference is negligible.
  \item UNC-3.R Determine whether a sampling distribution of a sample mean can be described as approximately normal. [Skill 3.C]
  \item UNC-3.R.1 For a numerical variable, if the population distribution can be modeled with a normal distribution, the sampling distribution of the sample mean can be modeled with a normal distribution. UNC-3.R.2 For a numerical variable, if the population distribution cannot be modeled with a normal distribution, the sampling distribution of the sample mean can be modeled approximately by a normal distribution provided the sample size is large enough (n $\geq$ 30).
  \item UNC-3.S Interpret probabilities and parameters for a sampling distribution for a sample mean. [Skill 4.B]
  \item UNC-3.S.1 Probabilities and parameters for a sampling distribution for a sample mean should be interpreted using appropriate units and within the context of a specific population.
\end{itemize}
}
\textbf{Essential question:} How do population shape, sample size, and sampling variability determine whether a normal probability for a sample mean is trustworthy?\par
\textbf{DOK-3 skill:} 4.B \textperiodcentered\ \textbf{FRQ pattern:} {\small\texttt{normal-\allowbreak{}vs-\allowbreak{}clt-\allowbreak{}model-\allowbreak{}adjudication}}\par
\textbf{DOK rationale:} Part (c) adjudicates two incomplete normal-model claims by coordinating population shape, sample size, the 10 percent condition, sampling-distribution variability, and the consequence of rejecting a valid probability model.\par

\textbf{Self-paced bonus sheet.} Students take it when they choose and turn it in when they finish. Everything they need is printed on the sheet. Score part (c) only, E / P / I, by hand --- never AI-graded or auto-scored. (a) and (b) are the ladder up; use them to see where a thin (c) came from.\par

\textbf{Questions to ask}
\begin{itemize}[leftmargin=1.5em,itemsep=2pt,topsep=2pt]
  \item Is 30 the only route to a normal sampling distribution?
  \item Does the Central Limit Theorem change the original population or the distribution of sample means?
  \item Why are both sampling distributions centered at 40 despite different sample sizes?
  \item How do standard deviations 4 and 2 change where 46 falls on each model?
\end{itemize}

\begin{callout}{\IconWarn\ LOOK FOR}{calloutred}
\begin{itemize}[leftmargin=1.5em,itemsep=2pt,topsep=2pt]
  \item Using n at least 30 as a universal requirement even for a normal population.
  \item Claiming a skewed population prevents the CLT from applying at n equals 36.
  \item Using the population standard deviation of 12 for both distributions of sample means.
\end{itemize}
\end{callout}

\sectionbanner{SCORING PART (c) --- E / P / I}

\textbf{Expected elements}
\begin{itemize}[leftmargin=1.5em,itemsep=2pt,topsep=2pt]
  \item \textbf{judgment-1} (required) --- Supports both approximations, using A population normality and B n=36 CLT route.
  \item \textbf{judgment-2} (required) --- Explains the correct and incorrect portion of both students' claims.
  \item \textbf{judgment-3} (required) --- Connects SDs 4 and 2, or z-scores 1.5 and 3, to the different tail probabilities.
\end{itemize}

{\small\renewcommand{\arraystretch}{1.25}
\noindent\begin{tabular}{|p{0.5in}|p{5.9in}|}
\hline
\textbf{E} & Meets all three checks with correct contextual reasoning; equivalent wording is acceptable. \\ \hline
\textbf{P} & Shows at least one substantive correct check, but has an omission or error that prevents E. \\ \hline
\textbf{I} & Shows none of the three checks with substantive correct reasoning; a choice alone is insufficient. \\ \hline
\end{tabular}}

\clearpage
\focusbanner{THE PROBLEM --- annotated}

Suppose Lines A and B each contain 1{,}000 cycle times with population mean 40 seconds and standard deviation 12 seconds. Line A is approximately normal and an SRS of $n_A=9$ is selected; Line B is strongly right-skewed and an SRS of $n_B=36$ is selected. An auditor wants $P(\bar X>46)$ for each.\par\smallskip \textbf{Serena:} ``Only B can use a normal model because $36\geq30$; its probability is $0.00135$. A is too small.''\par \textbf{Mateo:} ``Only A can use a normal model because its population is normal; its probability is $0.0668$. B is too skewed.''\par
\begin{center}
\textbf{Sampling plans}\par\smallskip
\begin{tabular}{|l|c|c|c|c|c|}
\hline
\textbf{Line} & \textbf{Shape} & \textbf{$N$} & \textbf{$n$} & \textbf{$\mu$} & \textbf{$\sigma$} \\ \hline
\textbf{A} & Normal & 1000 & 9 & 40 & 12 \\ \hline
\textbf{B} & Right-skewed & 1000 & 36 & 40 & 12 \\ \hline
\end{tabular}
\end{center}
\begin{firsttakebox}{FIRST TAKE --- one sentence before you work the parts}
Which plan seems more likely to give averages close to 40 seconds? Give one reason from the sample sizes or shapes.\par
\answer{Not graded. Each student is correct about one line and wrong about the other; the finish should replace either-or reasoning with the two separate normality pathways.}
\end{firsttakebox}

\begin{callout}{\IconBook\ MODEL SAMPLE MEANS, NOT INDIVIDUAL VALUES (keep on the page)}{calloutgreen}
For an SRS of size $n$, $\mu_{\bar X}=\mu$ and $\sigma_{\bar X}=\sigma/\sqrt n$ when
$n\leq0.10N$ for sampling without replacement. The sampling distribution of $\bar X$ is
approximately normal if the population is approximately normal (for any $n$)
\textbf{or} if the population is not normal but $n\geq30$ (Central Limit Theorem).
A normal tail probability describes the mean of an entire random sample, not one individual value.
For a normal probability, standardize with $z=(\text{value}-\text{mean})/\text{SD}$ and use the appropriate normal tail.
Printed normal table: $P(Z\leq1.5000)=0.933193$ and $P(Z\leq3.0000)=0.998650$. A right-tail area is 1 minus the left area.
\end{callout}

\dokbadge{1}~\textbf{(a)} State the mean of each sampling distribution and the 10\% condition for using $\sigma_{\bar X}=\sigma/\sqrt n$.\par
\answer{Both sampling distributions have mean 40 seconds. Without replacement, each sample must be no more than 10 percent of its population to use $\sigma_{\bar X}=\sigma/\sqrt n$ without a finite-population correction.}

\dokbadge{2}~\textbf{(b)} Verify the 10 percent condition, compute both sampling-distribution standard deviations, and calculate each normal tail probability when justified.\par
\answer{Both SRS sizes satisfy the cap $0.10(1{,}000)=100$. For A, $\sigma_{\bar X}=12/\sqrt9=4$; its normal population justifies $z=(46-40)/4=1.5$ and an upper tail of $0.0668$. For B, $\sigma_{\bar X}=12/\sqrt{36}=2$; the CLT applies because $36\geq30$, giving $z=3$ and an upper tail of $0.00135$.}

\dokbadge{3}~\textbf{(c)} In 3--4 sentences, explain what each student gets right and wrong. Use the two normality pathways to decide which probabilities are supported. Connect the different tail probabilities to the sampling SDs in (b).\par
\answer{Neither whole claim is complete. Serena is right about B: despite population skew, $n_B=36$ activates the CLT and supports $0.00135$; she wrongly rejects A, whose normal population supports a normal sampling distribution at $n_A=9$. Mateo is right about A and $0.0668$ but ignores B's CLT route. Their rules would each discard one valid probability. Both models center at 40, but 46 is 1.5 sampling SDs above A's center and 3 above B's, so B's larger sample produces the smaller tail.}

\end{document}
